\clearpage\fancyhead[L]{\small PALLAS.}\hypertarget{native-388-lower}{}\label{native-388-lower}
\begin{center}\textbf{PALLAS.}\par\vspace{4pt}\hrule\vspace{4pt}\fontsize{10}{12}\selectfont Astronomisches Jahrbuch für 1818. Seite 167\;..\;173. Berlin 1815.\par\vspace{4pt}\hrule\vspace{4pt}\end{center}

\begin{flushright}
Göttingen, den 1.\ Februar 1815.
\end{flushright}

\par\indent Die Berechnungen für \emph{Pallas}, die auch im letzten Jahrbuch stehen, haben sich erneut bestens bestätigt. Ich selbst habe den Planeten zwar nur einmal mit dem Kreismikrometer beobachtet

\begin{center}
\fontsize{9}{11}\selectfont\setlength{\tabcolsep}{3pt}\renewcommand{\arraystretch}{1.15}\begin{tabular}{l|ccc|ccc|ccc}
1814 & \multicolumn{3}{c|}{Mittlere Zeit} & \multicolumn{3}{c|}{Rektaszension} & \multicolumn{3}{c}{Deklination} \\
\hline
Sept.\ 16 & $12^{\mathrm u}$ & $1^{\mathrm m}$ & $1^{\mathrm s}$ & $146^\circ$ & $16'$ & $23''4$ & $11^\circ$ & $4'$ & $20''9$ S. \\
\end{tabular}
\end{center}

\noindent doch eine wertvolle Reihe von Meridianbeobachtungen der Herren von \textsc{Lindenau}, \textsc{Bessel} und \textsc{Schumacher} ermöglichte es mir, die Opposition gut zu berechnen. Für die Einzelheiten verweise ich auf die Gött.\ gel.\ Anz., St.\ 199 [S.\ 379 d.\ B.]. Hier führe ich nur das Endergebnis an:

\begin{center}25. Oktober 1814. $12^{\mathrm u}\;33^{\mathrm m}\;22^{\mathrm s}$ Mittlere Zeit in Göttingen\end{center}
\begin{center}\fontsize{10}{12}\selectfont\setlength{\tabcolsep}{3pt}\renewcommand{\arraystretch}{1.15}\begin{tabular}{lcr}
Wahre Länge & $.\;.\;.\;.\;.\;.\;.\;.\;.\;.\;.\;.\;.$ & $31^\circ\;58'\;11''3$\\
Geozentrische Breite, südl. & $.\;.\;.\;.\;.\;.\;.\;.$ & $37\;20\;53.2$\\
\end{tabular}\end{center}

\par\indent Meine Theorie gibt alle 10 bisherigen Oppositionen bis auf wenige Sekunden wieder.

\Needspace{12\baselineskip}\par\indent Herr \textsc{Enke} hat die verdienstvolle Arbeit übernommen, die Berechnungen für die nächste Erscheinung im Voraus anzustellen. Hier gebe ich sein Ergebnis für die XI.\ Opposition an:

\begin{center}9. April 1816. $2^{\mathrm u}\;14^{\mathrm m}\;24^{\mathrm s}$ Mittlere Zeit in Göttingen\end{center}
\begin{center}\fontsize{10}{12}\selectfont\setlength{\tabcolsep}{3pt}\renewcommand{\arraystretch}{1.15}\begin{tabular}{lcr}
Wahre Länge & $.\;.\;.\;.\;.\;.\;.\;.\;.\;.\;.\;.$ & $199^\circ\;34'\;40''5$\\
Geozentrische Breite & $.\;.\;.\;.\;.\;.\;.\;.\;.$ & $28\;6\;16.3\;\text{Nordl.}$\\
\end{tabular}\end{center}

\Needspace{8\baselineskip}\par\indent Bei der nächsten Opposition wird \emph{Pallas} heller sein, als sie seit mehreren Jahren gewesen ist. Hier die Lichtstärken zum Vergleich

\begin{center}\fontsize{9}{11}\selectfont
1812\ .\ 0.01666,\quad 1813\ .\ 0.01475,\quad 1814\ .\ 0.05476,\quad 1816\ .\ 0.05997
\end{center}

\par\indent Ich freue mich, Ihnen hier eine Abschrift der Ephemeride beilegen zu können, die derselbe geschickte junge Astronom berechnet hat.

\par\indent Für die Beobachtung der Wintersonnenwende war das Wetter in diesem Jahr ungewöhnlich günstig. Ich habe alles getan, um diese Beobachtungen so genau wie möglich zu machen. Dazu \SourceLeaf{389}gehörte ein eigener Schirm, der Licht nur auf das Objektiv fallen ließ und alle anderen Strahlen abhielt. Da ich diese Vorrichtung bei meinen Beobachtungen früherer Sonnenwenden noch nicht angebracht hatte, wollte ich diesen Beobachtungen kein volles Vertrauen schenken. Allerdings kann ich eigentlich nicht sagen, dass die Beobachtungen vor dem Anbringen des Schirms schlechter übereinstimmten. Das Ergebnis für die Schiefe der Ekliptik stimmt nun vollständig mit dem von Herrn \textsc{von Zach} überein. Die Wintersonnenwende ergibt einen um $10''$ kleineren Wert als die Sommersonnenwende. Ich möchte dieses Phänomen jedoch lieber so ausdrücken. Legt man für die Schiefe der Ekliptik den Mittelwert zugrunde, so ergeben Sonnenbeobachtungen stets eine im Durchschnitt um $5''$ geringere Polhöhe als Beobachtungen von Zirkumpolarsternen. Nach meiner Erfahrung gilt das zu \emph{jeder} Jahreszeit. Dieses sonderbare Phänomen ist noch immer ein astronomisches Rätsel. Es wäre gelöst, wenn sich am Instrument etwas fände, das die Annahme rechtfertigte, es gebe alle Zenitdistanzen um $2''5$ zu klein an. Ich kann aber überhaupt nichts finden, was stets einen Fehler in derselben Richtung hervorriefe. Besonders möchte ich ausdrücklich darauf hinweisen, dass ich die Visierlinie des Fernrohrs mit einem eigenen Verfahren so genau wie möglich parallel zur Ebene des Instruments ausgerichtet habe. Dass ich zugleich darauf geachtet habe, die Ebene genau senkrecht zu stellen und bei der Berechnung nichts auszulassen, brauche ich nicht eigens zu erwähnen. So wurden beispielsweise die von der vierten Potenz abhängigen Glieder jedes Mal berücksichtigt. Wenn man auch nur eine dieser Vorsichtsregeln außer Acht lässt, wird die Zenitdistanz tatsächlich zu klein. Meine Beobachtungen der letzten Sonnenwende waren folgende:

\begin{center}
\fontsize{10}{12}\selectfont\setlength{\tabcolsep}{5pt}\renewcommand{\arraystretch}{1.15}\begin{tabular}{l l ccc l}
1814 Dez.\ 17 & W.\ Z.\ D.\ & $74^\circ$ & $53'$ & $48''53$ & 12 Beob. \\
19 & & 74 & 57 & 35.13 & 14 \\
21 & & 74 & 59 & 24.02 & 14 \\
22 & & 74 & 59 & 38.63 & 14 \\
23 & & 74 & 59 & 21.10 & 14 \\
24 & & 74 & 58 & 36.74 & 14 \\
25 & & 74 & 57 & 22.29 & 14 \\
29 & & 74 & 57 & 51.39 & 12 \\
\end{tabular}
\end{center}

\par\indent An jedem Tag waren die Beobachtungen in 2 oder 3 Reihen aufgeteilt. Die einzelnen Reihen wichen meist kaum eine und nie zwei Sekunden vom Mittel ab.

\par\indent Für eine ganz genaue Berechnung dieser Beobachtungen hat mir bisher die Zeit gefehlt. Vorläufig ergibt sich daraus aber für die Zenitdistanz bei der Sonnenwende $74^\circ\,59'\,36''6$. Mit der Polhöhe $51^\circ\,31'\,55''0$ ergibt das eine Schiefe, die um $5''$ kleiner ist als in unseren besten Tafeln. Die Polhöhe beruht auf mehr als 200 Beobachtungen des Polarsterns bei seiner oberen und unteren Kulmination. Die untere ergab $51^\circ\,31'\,55''5$, die obere $51^\circ\,31'\,54''5$, wenn ich die Tafeln von Herrn \textsc{Bessel} zugrunde legte.

\Needspace{17\baselineskip}\begin{flushright}
Göttingen, den 24.\ April
\end{flushright}

\par\indent Den Kometen von \textsc{Olbers} werden Sie inzwischen zweifellos längst selbst gefunden und beobachtet haben. Ich fand ihn zwar gleich am 13. März, als ich die erste Nachricht von seiner Entdeckung erhielt. Das äußerst ungünstige Wetter ließ aber erst am 20. eine wirkliche Beobachtung zu. Meine Beobachtungen aus dem März sind folgende:

\begin{center}\fontsize{9}{11}\selectfont\setlength{\tabcolsep}{3pt}\renewcommand{\arraystretch}{1.15}\begin{tabular}{l|ccc|ccc|ccc}
1815 & \multicolumn{3}{c|}{Mittlere Zeit} & \multicolumn{3}{c|}{\shortstack{Scheinbare\\Rektaszension}} & \multicolumn{3}{c}{\shortstack{Scheinbare Dekl.\\Nördlich}} \\

\hline
März 20 & $10^{\mathrm u}$ & $33^{\mathrm m}$ & $6^{\mathrm s}$ & $54^\circ$ & $7'$ & $1''$ & $39^\circ$ & $7'$ & $47''$ \\
21 & 10 & 11 & 37 & 54 & 33 & 21 & 39 & 36 & 57 \\
21 & 10 & 57 & 0 & 54 & 34 & 21 & 39 & 38 & 8 \\
25 & 9 & 46 & 38 & 56 & 28 & 50 & 41 & 38 & 5 \\
30 & 9 & 50 & 57 & 59 & 14 & 3 & 44 & 10 & 27 \\
\end{tabular}\end{center}

\Needspace{20\baselineskip}\SourceLeaf{390}\par\indent Die zweite Beobachtung vom 21. stammt von Herrn \textsc{Enke}, der auch sämtliche Vergleichssterne aus der \emph{Histoire Céleste} reduziert hat. Meine Beobachtungen vom April sind noch nicht reduziert. --- Am 31. März berechnete ich aus den obigen Beobachtungen die folgenden vorläufigen parabolischen Elemente, die auch jetzt noch nicht stark abweichen dürften.

\begin{center}\fontsize{9}{11}\selectfont $\text{Perihel}\left\{\begin{array}{lcr}
\text{Zeit, 24. April 1815} & .\;.\;.\;. & 16^{\mathrm u}\;37^{\mathrm m}\;34^{\mathrm s}\;\text{Göttinger Zeit}\\
\text{Länge} & .\;.\;.\;.\;.\;.\;.\;.\;. & 146^\circ\;7\;2\\
\text{Abstand} & .\;.\;.\;.\;.\;.\;.\;. & 1.24738\\
\end{array}\right.$\end{center}
\begin{center}\fontsize{10}{12}\selectfont\setlength{\tabcolsep}{3pt}\renewcommand{\arraystretch}{1.15}\begin{tabular}{lcr}
Aufsteigender Knoten & $.\;.\;.\;.\;.\;.\;.\;.$ & $82\;43\;4$\\
Neigung der Bahn & $.\;.\;.\;.\;.\;.\;.\;.\;.$ & $45\;8\;55$\\
\end{tabular}\end{center}
\noindent Bewegung rechtläufig.

\par\indent Ich hoffe, dass sich dieser Komet noch bis Juli beobachten lässt und seine Bahn anschließend mit großer Zuverlässigkeit als Ellipse berechnet werden kann.

\par\indent Juno war um die Zeit ihrer Opposition nur von 9. bis 10. Größe. Ich habe sie zweimal gut beobachtet, die Beobachtungen aber noch nicht reduziert.

\par\indent Unser ausgezeichnetes Heliometer konnte bisher mangels einer festen Aufstellung nur wenig genutzt werden. Vorgestern ist endlich das parallaktische Stativ angekommen, das seit fast einem Jahr bestellt war und von \textsc{Frauenhofer} äußerst schön und einfallsreich ausgeführt wurde.
